En un jaciment arqueològic es troben unes restes òssies antigues d'animals. Un gram d'aquestes restes conté $9,5\times 10^{8}$ àtoms de carboni 14. L'anàlisi d'una mostra actual, de la mateixa massa i de característiques similars, revela que, en el moment de la mort dels animals, els ossos tenien $6,9\times 10^{9}$ àtoms de C-14/gram.

\begin{apartat}{1}
Determineu l'antiguitat de les restes si sabem que el període de semidesintegració del C-14 és de 5\,760 anys.
\end{apartat}

\begin{apartat}{1}
Escriviu l'equació nuclear de la desintegració (amb emissió de $\beta^-$) del C-14 i incloeu-hi els antineutrins. Calculeu el defecte de massa per nucleó de C-14.
\end{apartat}

\dades{%
  $c = 3,00\times 10^{8}\ \mathrm{m\,s^{-1}}$\\
  Nombres atòmics: Be, 4; B, 5; C, 6; N, 7; O, 8; F, 9}

\noindent Masses:
\begin{center}
\begin{tabular}{|l|r|l|r|}
\hline
\textit{Partícula} & \textit{Massa (kg)} & \textit{Partícula} & \textit{Massa (kg)}\\
\hline
protó & $1,6726\times 10^{-27}$ & electró & $9,1093\times 10^{-31}$\\
\hline
neutró & $1,6749\times 10^{-27}$ & àtom de C-14 & $2,3253\times 10^{-26}$\\
\hline
\end{tabular}
\end{center}
